\pagestyle{empty}
\tight
\begin{tblr}{width=\textwidth,colspec={|Q[c,m,1.9cm]|X[l,m]|},hlines,vlines,hspan=minimal,row{1}={bg=headblue},rowsep=1pt,colsep=3pt}
\SetCell{c,m}\hfil\logo\hfil & \textbf{POLITEKNIK NEGERI MALANG}\par\textbf{JURUSAN TEKNOLOGI INFORMASI}\par\textbf{PROGRAM STUDI: D4 TEKNIK INFORMATIKA} \\
\SetCell[c=2]{c} \textbf{RENCANA PEMBELAJARAN SEMESTER (RPS)} & \\
\end{tblr}
\begin{longtblr}{width=\textwidth,colspec={|Q[l,m,1.9cm]|X[0.85,l,m]|X[1.45,l,m]|X[0.75,l,m]|X[1.15,l,m]|},hlines,vlines,hspan=minimal,rowsep=1pt,colsep=2pt,row{1,3}={bg=headgray,font=\bfseries}}
MATA KULIAH & KODE & BOBOT (SKS)/jam & SEMESTER & TGL. PENYUSUNAN \\
Praktikum Algoritma dan Struktur Data & RTI252009 & 3 SKS/6 jam & 2 & 31 Januari 2026 \\
OTORISASI & Dosen Pengembang RPS & Koordinator RMK & \SetCell[c=2]{l} Ka PRODI &  \\
 & Mungki Astiningrum, S.T., M.Kom.\par Mamluatul Hani'ah, S.Kom., M.Kom.\par Vivi Nur Wijayaningrum, S.Kom., M.Kom. & Mungki Astiningrum, S.T., M.Kom. & \SetCell[c=2]{l}{\vspace{8mm}\par Dr. Ely Setyo Astuti, S.T., M.T.} &  \\
\SetCell[r=11]{l} Capaian Pembelajaran (CP) & \SetCell[c=4]{l,bg=headgray,font=\bfseries} Capaian Pembelajaran Lulusan Program Studi (CPL-Prodi) &  &  &  \\
 & CPL07 & \SetCell[c=3]{l} Mampu mengembangkan sistem komputasi cerdas, analisis data dengan menerapkan algoritma dan teknik kecerdasan artifisial secara tepat guna. &  &  \\
 & CPL09 & \SetCell[c=3]{l} Memiliki pemahaman yang memadai terkait cara kerja sistem komputer dan berbagai teknik/pendekatan berbasis komputasi untuk memecahkan masalah stakeholder. &  &  \\
 & \SetCell[c=4]{l,bg=headgray,font=\bfseries} Capaian Pembelajaran mata kuliah (CPL-MK) &  &  &  \\
 & CPMK0703 & \SetCell[c=3]{l} Mampu memilih dan mengimplementasikan struktur data serta algoritma yang efisien untuk mendukung pengolahan data dan pengembangan sistem. &  &  \\
 & CPMK0708 & \SetCell[c=3]{l} Mampu menerapkan konsep struktur data dalam perancangan dan implementasi algoritma untuk menghasilkan solusi komputasi yang efisien. &  &  \\
 & CPMK0903 & \SetCell[c=3]{l} Mahasiswa mampu menjelaskan prinsip logika komputasi serta menerapkan teknik dasar pemrograman dan algoritma untuk menganalisis dan menyelesaikan permasalahan. &  &  \\
 & \SetCell[c=4]{l,bg=headgray,font=\bfseries} Kemampuan akhir tiap tahapan belajar (Sub-CPMK) &  &  &  \\
 & SCPMK0903-01801 & \SetCell[c=3]{l} Mahasiswa mampu menerapkan teknik dasar pemrograman dan algoritma untuk menyelesaikan permasalahan komputasi menggunakan array of object. &  &  \\
 & SCPMK0708-01801 & \SetCell[c=3]{l} Mahasiswa mampu mengimplementasikan struktur data linear dan mengintegrasikannya ke dalam algoritma untuk menyelesaikan permasalahan komputasi. &  &  \\
 & SCPMK0703-01801 & \SetCell[c=3]{l} Mahasiswa mampu memilih dan mengimplementasikan struktur data non-linear serta koleksi data untuk membangun solusi pengolahan data yang efisien. &  &  \\
\SetCell[r=2]{l} Penilaian & \SetCell[c=4]{l} *Nilai CPMK didapat dari agregasi nilai SUB-CPMK yang dibebankan &  &  &  \\
 & \SetCell[c=4]{l}{\tiny\begin{tblr}{width=\linewidth,colspec={|Q[c,m,0.09\linewidth]|Q[c,m,0.15\linewidth]|X[l,m]|X[l,m]|X[l,m]|X[l,m]|X[l,m]|Q[c,m,0.08\linewidth]|},hlines,vlines,hspan=minimal,rowsep=1pt,colsep=0.7pt,row{1,2}={bg=headgray,font=\bfseries}}
Kode CPMK & Kode SCPMK & \SetCell[c=5]{c} Bobot per Bentuk Penilaian & & & & & Total Bobot \\
 & & Tugas & Tes Praktik 1\par(Asessment 1) & Tes Case Method 1\par(Asessment 2) & Tes Case Method 2\par(Asessment 1) & Tes Praktik 2\par(Asessment 1) & \\
CPMK0903 & SCPMK0903-01801 & 5\%: materi 1,2,3,5,6,7 & 20\%: materi 1--3 & 25.5\%: materi 5--7 & & & 50.5\% \\
CPMK0708 & SCPMK0708-01801 & 2\%: materi 9,10,11,12 & & & 25.5\%: materi 9--12 & & 27.5\% \\
CPMK0703 & SCPMK0703-01801 & 2\%: materi 14,15 & & & & 20\%: materi 14--15 & 22\% \\
\SetCell[c=7]{r}\textbf{Total Per Penilaian} & & & & & & & \textbf{100\%} \\
\end{tblr}} &  &  &  \\
Diskripsi Singkat Mata Kuliah & \SetCell[c=4]{l} Matakuliah Praktikum Algoritma dan Struktur Data memberikan pengetahuan dan keterampilan praktis pembuatan algoritma dan struktur data yang meliputi brute force, divide-conquer, stack, queue, linked list, tree, dan graph serta proses sorting dan searching. &  &  &  \\
Bahan Kajian / Materi Pembelajaran & \SetCell[c=4]{l}\begin{enumerate}\item Searching\item Sorting\item Queue\item Stack\item Linked List\item Tree\item Graph\item Brute Force\item Divide-Conquer\item Depth First Search\item Breadth First Search\end{enumerate} &  &  &  \\
\SetCell[r=2]{l} Pustaka & \SetCell[c=4]{l} \textbf{Utama:}\begin{enumerate}\item Goodrich, M. T., Tamassia, R., and Goldwasser, M. H. 2014. Data Structures \& Algorithms in Java, 6th Edition. Wiley Global Education.\item Ramadhani, C. 2015. Dasar Algoritma dan Struktur Data dengan Bahasa Java. Yogyakarta: Andi Publisher.\item Nugroho, A. 2008. Algoritma dan Struktur Data dalam Bahasa Java. Yogyakarta: Andi Publisher.\item Hariyanto, B. 2007. Struktur Data. Bandung: Informatika.\end{enumerate} &  &  &  \\
 & \SetCell[c=4]{l} \textbf{Pendukung:}\begin{enumerate}\item Buana, I. S., Nata, G. N. M., and Arnawa, I. K. 2018. Struktur Data. Yogyakarta: Andi Publisher.\item Kadir, A. 2015. Teori dan Aplikasi Struktur Data Menggunakan Java. Yogyakarta: Andi Publisher.\end{enumerate} &  &  &  \\
Media Pembelajaran & \SetCell[c=2]{l} \textbf{Software:} IDE Java, compiler/JDK, aplikasi flowchart, dan LMS (lms.polinema.ac.id). &  & \SetCell[c=2]{l} \textbf{Hardware:} Komputer laboratorium, laptop, dan proyektor. &  \\
Nama Dosen Pengampu & \SetCell[c=4]{l}\begin{enumerate}\item Mamluatul Hani'ah, S.Kom., M.Kom.\item Vivi Nur Wijayaningrum, S.Kom., M.Kom.\item Imam Fahrur Rozi, S.T., M.T.\item Novta Dany'el Irawan, S.Pd., M.T.\item Mungki Astiningrum, S.T., M.Kom.\item Irsyad Arif Mashudi, S.Kom., M.Kom.\item Rokhimatul Wakhidah, S.Pd., M.T.\end{enumerate} &  &  &  \\
Matakuliah Syarat & \SetCell[c=4]{l} RTI251006 Dasar Pemrograman &  &  &  \\
\end{longtblr}
\begin{longtblr}{width=\textwidth,colspec={|Q[c,m,0.06\textwidth]|X[1.35,l,m]|X[1.08,l,m]|X[1.16,l,m]|X[0.7,l,m]|X[1.24,l,m]|X[1.06,l,m]|X[0.98,l,m]|Q[c,m,0.05\textwidth]|},hlines,vlines,hspan=minimal,rowhead=2,row{1,2}={bg=headgreen,font=\bfseries},rowsep=1pt,colsep=1.5pt}
Mgg.\par Ke & Kemampuan Akhir Yang Direncanakan (Sub-CP-MK) & Bahan kajian (Materi Pembelajaran) & Bentuk dan Metode Pembelajaran & Estimasi Waktu & Pengalaman Belajar Mahasiswa & Kriteria \& Bentuk Penilaian & Indikator Penilaian & Bobot Penilaian (\%) \\
(1) & (2) & (3) & (4) & (5) & (6) & (7) & (8) & (9) \\
1 & SCPMK0903-01801 & Pemilihan, perulangan, array, dan fungsi. & Modalitas: Blended Learning. Bentuk: Luring. Metode: diskusi kelompok dan studi kasus. Sumber belajar: LMS Polinema. & 1 x 3 x 50' praktikum (asinkron video, sinkron tatap muka, tugas terstruktur); penugasan 1 x 50'. & Mahasiswa memahami konsep dasar algoritma, mengimplementasikan flowchart untuk studi kasus (deret bilangan, membalik kata pada array char, konversi detik), dan mempresentasikan hasil. & Rubrik kriteria penilaian; keaktifan diskusi; tes tulis atau tugas pemecahan studi kasus dengan flowchart. & Pemahaman konsep dasar algoritma; ketepatan pembuatan flowchart untuk menyelesaikan studi kasus. & 1 \\
2 & SCPMK0903-01801 & Konsep objek, class, atribut, method, instansiasi, konstruktor, dan class diagram. & Modalitas: Blended Learning. Bentuk: Luring. Metode: diskusi kelompok dan studi kasus. Sumber belajar: LMS Polinema. & 1 x 3 x 50' praktikum (asinkron video, sinkron tatap muka, tugas terstruktur); penugasan 1 x 50'. & Mahasiswa memahami object, mendeklarasikan class, atribut, method, membuat object, mengakses atribut dan method, memanggil konstruktor, serta menuangkan object ke class diagram. & Rubrik kriteria penilaian; keaktifan diskusi; tes tulis atau tugas pembuatan class diagram. & Pemahaman konsep object, class, atribut, dan method; ketepatan pembuatan class diagram. & 1 \\
3 & SCPMK0903-01801 & Array of object: deklarasi, instansiasi, assignment, dan penampilan data dengan bahasa Java. & Modalitas: Blended Learning. Bentuk: Kuliah luring. Metode: diskusi kelompok virtual dan studi kasus. Sumber belajar: LMS Polinema. & 1 x 3 x 50' praktikum (asinkron video, sinkron tatap muka, tugas terstruktur); penugasan 1 x 50'. & Mahasiswa memahami konsep array of object dan menerapkan array of object dalam flowchart untuk berbagai model studi kasus. & Rubrik kriteria penilaian; keaktifan diskusi; tes tulis atau tugas pemecahan studi kasus dengan flowchart. & Pemahaman pembuatan array of object; ketepatan flowchart sebagai implementasi algoritma untuk menyelesaikan studi kasus. & 1 \\
4 & SCPMK0903-01801 & Tes Praktik 1 (Asessment 1): materi pekan 1--3. & Ujian praktik di laboratorium, mengerjakan program secara mandiri. & 1 x 3 x 50' sesuai jadwal asesmen. & Mahasiswa mengimplementasikan program berbasis array of object untuk studi kasus secara mandiri. & Ujian praktik; rubrik penilaian RTM. & Program terkompilasi dan berjalan benar; ketepatan implementasi class dan array of object; kualitas kode dan penjelasan. & 20 \\
5 & SCPMK0903-01801 & Algoritma brute force, divide-conquer, kompleksitas Big-O, dan perhitungan notasi Big-O. & Modalitas: Blended Learning. Bentuk: Kuliah luring. Metode: diskusi kelompok virtual dan studi kasus. Sumber belajar: LMS Polinema. & 1 x 3 x 50' praktikum (asinkron video, sinkron tatap muka, tugas terstruktur); penugasan 1 x 50'. & Mahasiswa memahami konsep brute force dan divide-conquer serta mengimplementasikannya dalam flowchart untuk berbagai model studi kasus. & Rubrik kriteria penilaian; keaktifan diskusi; tes tulis atau tugas pemecahan studi kasus dengan flowchart. & Pemahaman brute force dan divide-conquer; ketepatan flowchart sebagai implementasi algoritma. & 1 \\
6 & SCPMK0903-01801 & Sorting: bubble sort, selection sort, dan insertion sort. & Modalitas: Blended Learning. Bentuk: Kuliah luring. Metode: diskusi kelompok virtual dan studi kasus. Sumber belajar: LMS Polinema. & 1 x 3 x 50' praktikum (asinkron video, sinkron tatap muka, tugas terstruktur); penugasan 1 x 50'. & Mahasiswa memahami konsep algoritma sorting, menentukan algoritma yang digunakan, dan membuat flowchart untuk studi kasus sorting. & Rubrik kriteria penilaian; keaktifan diskusi; tes tulis atau tugas pemecahan studi kasus dengan flowchart. & Pemahaman algoritma sorting; ketepatan flowchart sebagai implementasi algoritma untuk menyelesaikan studi kasus. & 1 \\
7 & SCPMK0903-01801 & Sequential search, binary search, dan pengayaan merge sort. & Modalitas: Blended Learning. Bentuk: Luring. Metode: diskusi kelompok dan studi kasus. Sumber belajar: LMS Polinema. & 1 x 3 x 50' praktikum (asinkron video, sinkron tatap muka, tugas terstruktur); penugasan 1 x 50'. & Mahasiswa memahami konsep searching, membuat pseudocode sequential search dan binary search, dan melakukan simulasi manual penyelesaian kasus pencarian. & Ketepatan dan penguasaan; keaktifan diskusi; ketepatan jawaban tugas minggu 7. & Pemahaman penerapan searching; kesesuaian jawaban tugas. & 0 \\
8 & SCPMK0903-01801 & Tes Case Method 1 (Asessment 2): materi pekan 5--7. & Ujian berbasis case method, mengerjakan penyelesaian kasus algoritmik secara mandiri. & 1 x 3 x 50' sesuai jadwal asesmen. & Mahasiswa menyelesaikan kasus sorting dan searching dengan memilih algoritma yang sesuai dan mengimplementasikannya. & Ujian case method; rubrik penilaian RTM. & Ketepatan analisis kasus dan pemilihan algoritma; kebenaran implementasi pada kasus uji; kualitas dokumentasi. & 25.5 \\
9 & SCPMK0708-01801 & Konsep struktur data stack, operasi-operasi pada stack, dan postfix expressions. & Modalitas: Blended Learning. Bentuk: Luring. Metode: Contextual Teaching and Learning (CTL). Sumber belajar: LMS Polinema. & 1 x 3 x 50' praktikum (asinkron video, sinkron tatap muka, tugas terstruktur); penugasan 1 x 50'. & Mahasiswa memahami struktur data stack, operasi pada stack, dan mengimplementasikan konsep stack untuk menyelesaikan studi kasus. & Ketepatan dan penguasaan; keaktifan diskusi; ketepatan jawaban tugas minggu 9. & Pemahaman penerapan konsep stack; kesesuaian jawaban tugas. & 0.5 \\
10 & SCPMK0708-01801 & Konsep struktur data queue dan operasi-operasi pada queue. & Modalitas: Blended Learning. Bentuk: Luring. Metode: Contextual Teaching and Learning (CTL). Sumber belajar: LMS Polinema. & 1 x 3 x 50' praktikum (asinkron video, sinkron tatap muka, tugas terstruktur); penugasan 1 x 50'. & Mahasiswa memahami struktur data queue, operasi pada queue, dan mengimplementasikan konsep queue untuk menyelesaikan studi kasus. & Ketepatan dan penguasaan; keaktifan diskusi; ketepatan jawaban tugas minggu 10. & Pemahaman penerapan konsep queue; kesesuaian jawaban tugas. & 0.5 \\
11 & SCPMK0708-01801 & Linked list: insert data pada awal node, sebelum/sesudah sebuah node, akhir node, dan penghapusan node. & Modalitas: Blended Learning. Bentuk: Kuliah. Metode: diskusi kelompok dan studi kasus. Sumber belajar: LMS Polinema. & 1 x 3 x 50' praktikum (asinkron video, sinkron tatap muka, tugas terstruktur); penugasan 1 x 50'. & Mahasiswa menerapkan konsep linked list dalam diagram flowchart dan menunjukkan hasil pemahaman konsep linked list. & Rubrik kriteria penilaian; presentasi; tes tulis atau tugas pemecahan studi kasus dengan flowchart. & Pemahaman konsep linked list; ketepatan diagram flowchart sesuai studi kasus; performa presentasi. & 0.5 \\
12 & SCPMK0708-01801 & Double linked list. & Modalitas: Blended Learning. Bentuk: Kuliah. Metode: diskusi kelompok dan studi kasus. Sumber belajar: LMS Polinema. & 1 x 3 x 50' praktikum (asinkron video, sinkron tatap muka, tugas terstruktur); penugasan 1 x 50'. & Mahasiswa memahami double linked list sebagai pengembangan linked list, menerapkannya dalam diagram flowchart, dan menunjukkan hasil pemahaman. & Rubrik kriteria penilaian; presentasi; tes tulis atau tugas pemecahan studi kasus dengan flowchart. & Pemahaman konsep double linked list; ketepatan diagram flowchart sesuai studi kasus; performa presentasi. & 0.5 \\
13 & SCPMK0708-01801 & Tes Case Method 2 (Asessment 1): materi pekan 9--12. & Ujian berbasis case method, mengerjakan penyelesaian kasus struktur data linear secara mandiri. & 1 x 3 x 50' sesuai jadwal asesmen. & Mahasiswa menyelesaikan kasus dengan memilih dan mengimplementasikan struktur data linear (stack, queue, linked list) yang sesuai. & Ujian case method; rubrik penilaian RTM. & Ketepatan pemilihan struktur data linear; kebenaran implementasi operasi; kualitas demo dan penjelasan kode. & 25.5 \\
14 & SCPMK0703-01801 & Tree dan binary search tree. & Modalitas: Blended Learning. Bentuk: Kuliah. Metode: diskusi kelompok dan studi kasus. Sumber belajar: LMS Polinema. & 1 x 3 x 50' praktikum (asinkron video, sinkron tatap muka, tugas terstruktur); penugasan 1 x 50'. & Mahasiswa memahami konsep tree, penerapan binary tree, konsep binary search tree, dan tahapan penerapannya melalui studi kasus. & Ketepatan dan penguasaan; keaktifan diskusi; ketepatan jawaban tugas minggu 14. & Pemahaman konsep tree dan binary search tree; ketepatan menjawab studi kasus. & 1 \\
15 & SCPMK0703-01801 & Graph, konversi matriks ke graph, dan graph ke matriks. & Modalitas: Blended Learning. Bentuk: Kuliah. Metode: diskusi kelompok dan studi kasus. Sumber belajar: LMS Polinema. & 1 x 3 x 50' praktikum (asinkron video, sinkron tatap muka, tugas terstruktur); penugasan 1 x 50'. & Mahasiswa membuat algoritma graph secara umum dan menerapkan algoritma graph pada program. & Ketepatan dan penguasaan; keaktifan diskusi; ketepatan jawaban tugas minggu 15. & Pemahaman konsep graph; ketepatan menjawab tugas. & 1 \\
16 & SCPMK0703-01801 & Tes Praktik 2 (Asessment 1): materi pekan 14--15. & Ujian praktik di laboratorium, mengerjakan program secara mandiri. & 1 x 3 x 50' sesuai jadwal asesmen. & Mahasiswa mengimplementasikan struktur data non-linear (tree, binary search tree, graph) untuk studi kasus secara mandiri. & Ujian praktik; rubrik penilaian RTM. & Kebenaran implementasi tree/BST dan graph; kebenaran hasil pada kasus uji; kualitas penjelasan kode. & 20 \\
17 & SCPMK0903-01801, SCPMK0708-01801, SCPMK0703-01801 & Pengayaan materi pekan 1--15. & Pengayaan untuk mahasiswa sesuai Sub-CPMK yang belum tercapai. & Sesuai kebutuhan. & Mahasiswa mengerjakan pengayaan materi pekan 1--15 secara mandiri. & Ujian atau pengayaan. & Ketepatan jawaban dengan kunci jawaban. & 0 \\
\end{longtblr}
